%début contraintes ensae
\usepackage{pdfpages, setspace, mathptmx} %times roman
\usepackage{textpos}%pour textblock

% Début JMS
\makeatletter
\def\@maketitle{%
\thispagestyle{fancy}
\begin{center}
\includegraphics[width=15cm]{img/head_jms2022.png} 
\line(1,0){450}
\vspace{5mm}
\textbf{{\huge E}\Large stimation en temps réel de la tendance-cycle : apport de l’utilisation des filtres asymétriques dans la détection des points de retournement}
\end{center}

\begin{center}
\textit{Alain QUARTIER-LA-TENTE} \\
\vspace{2mm}
\textit{(*) Insee,  Département des études économiques (DEE), et  Laboratoire d'Économie et de Management de Nantes-Atlantique (LEMNA)}\\
\vspace{2mm}
\url{alain.quartier-la-tente@insee.fr} 
\end{center}

\vspace{5mm}
\small{{\bf Mots-clés.} séries temporelles,  tendance-cycle,  désaisonnalisation,  \emph{nowcasting}.}

\vspace{5mm}
\small{{\bf Domaines.} 10.2 Désaisonnalisation --- 10.3 Nowcasting.}

\begin{center}
\line(1,0){450}
\end{center}

\normalsize
}
\usepackage{fancyhdr} 
\pagestyle{fancy}

\renewcommand{\headrulewidth}{0pt}
\fancyhead[C]{} 
\fancyhead[L]{}
\fancyhead[R]{}

\renewcommand{\footrulewidth}{0pt}
\fancyfoot[C]{\footnotesize{$14^{es}$ Journées de méthodologie statistique de l'Insee (JMS) / Mars 2022 / PARIS}}
\fancyfoot[R]{\thepage}
%FIN JMS

\makeatother% cinsérer page de garde

\usepackage{stmaryrd}
\usepackage{multicol}
\usepackage{graphicx}
\usepackage{animate, dsfont, here, xspace}
%\usepackage{tikz}       
\usepackage{tikz,pgfplots}
 \pgfplotsset{compat=1.17}
%\includepdf[fitpaper=true, pages=-]{img/pdg.pdf}


\DeclareMathOperator{\e}{e}
\DeclareMathOperator{\Determinant}{det}
\renewcommand{\P}{\mathds{P}} %Apparement \P existe déjà ?
\newcommand\N{\mathds{N}}
\newcommand\Norm{\mathcal{N}}
\newcommand\R{\mathds{R}}
\newcommand\Z{\mathds{Z}}
\newcommand{\determinant}[1]{\Determinant\left(#1\right)}
%\newcommand\C{\mathds{C}}
%\newcommand\Z{\mathds{Z}}


\newcommand\1{\mathds{1}}
\newcommand{\E}[2][]{{\mathds{E}}_{#1}
  \def\temp{#2}\ifx\temp\empty
  \else
    \left[#2\right]
  \fi
}
\newcommand{\V}[2][]{{\mathds{V}}_{#1}
  \def\temp{#2}\ifx\temp\empty
  \else
    \left[#2\right]
  \fi
}
\newcommand\ud{\,\mathrm{d}}
\newcommand{\ps}[2]{\left\langle #1 \,,\, #2 \right\rangle}

% blocks
\usepackage{environ}
\usepackage[tikz]{bclogo}

\tikzstyle{titlestyle} =[draw=black!80,fill=black!20, text=black,
 right=10pt, rounded corners]
\mdfdefinestyle{symmaryboxstyle}{
	linecolor=black!80, backgroundcolor = black!5,
	skipabove=\baselineskip, innertopmargin=\baselineskip,
	innerbottommargin=\baselineskip,
	userdefinedwidth=\textwidth,
	middlelinewidth=1.2pt, roundcorner=5pt,
	skipabove={\dimexpr0.5\baselineskip+\topskip\relax},
	frametitleaboveskip=\dimexpr-\ht\strutbox\relax,
	innerlinewidth=0pt,
}
\newcounter{summarybox}
\NewEnviron{summary_box}[2][true]{%
\refstepcounter{summarybox} % on incrémente le compteur
\begin{mdframed}[style=symmaryboxstyle,
nobreak=#1,
frametitle={%
      \tikz[baseline=(current bounding box.east),outer sep=0pt]
      \node[titlestyle,anchor=east]
    {Encadré \thesummarybox{} --- #2};}
]
\vspace{-0.5em}
\BODY
\end{mdframed}
}

\usepackage{amsthm}
%\theoremstyle{remark}
\newtheorem*{remarque}{Remarque}

\usepackage{mathrsfs}
\usepackage{fontawesome5}

\usepackage[style=authoryear,uniquename=false, uniquelist=false]{biblatex}
\DefineBibliographyStrings{french}{andothers={et\addabbrvspace alii}}

